\section{Linux File System Internals}
\label{sec:vfs}

% Figure palette for this section, after the slides' kernel-structure diagrams
\definecolor{vpind}{HTML}{353598}   % struct files_struct
\definecolor{vpyel}{HTML}{FFFF66}   % NTFS functions
\tikzset{
  vpts/.style={draw=white, very thick, rounded corners=6pt, fill=datac},
  vpfs/.style={draw=white, very thick, rounded corners=5pt, fill=vpind},
  vpw/.style={draw=white, thick, text=white, font=\scriptsize, inner sep=2pt, minimum height=4.5mm},
  vpslot/.style={draw=white, fill=dirc, font=\scriptsize, minimum size=4.5mm, inner sep=0pt},
  vpsq/.style={draw=white, fill=dirc, minimum size=4mm, inner sep=0pt},
  vpmod/.style={fill=metac, font=\scriptsize\bfseries, align=center, inner sep=3pt},
  vpkern/.style={fill=kspace, rounded corners=3pt, inner sep=8pt,
                 label={[font=\small\itshape, anchor=south west, xshift=3pt, yshift=3pt]south west:Kernel memory}},
  vpdot/.style={-{Stealth[length=1.6mm]}, densely dotted, thick},
  % a stack of platters with "Disk" on top; anchors -top, -bot, -west, -east
  pics/vpdisk/.style={code={
    \foreach \y in {-0.6,-0.45,-0.3,-0.15,0} {
      \fill[black!45, draw=white, line width=0.8pt] (-0.62,\y+0.15) -- (-0.62,\y)
        arc[start angle=180, end angle=360, x radius=0.62, y radius=0.16] -- (0.62,\y+0.15) -- cycle;
    }
    \fill[black!45, draw=white, line width=0.8pt] (0,0.15) ellipse[x radius=0.62, y radius=0.16];
    \node[text=white, font=\scriptsize\bfseries] at (0,0.15) {Disk};
    \coordinate (-top) at (0,0.31); \coordinate (-bot) at (0,-0.76);
    \coordinate (-west) at (-0.62,-0.2); \coordinate (-east) at (0.62,-0.2);
  }},
  % a file: a page with a few text lines; anchors -west, -south
  pics/vpdoc/.style={code={
    \draw[fill=freec] (-0.4,-0.45) rectangle (0.4,0.45);
    \foreach \a/\b/\y in {-0.28/0.2/0.3, -0.1/0.3/0.13, -0.28/0.2/-0.04, -0.18/0.3/-0.25}
      \draw[thin] (\a,\y) -- (\b,\y);
    \coordinate (-west) at (-0.4,0); \coordinate (-south) at (0,-0.45);
  }},
  % \vptab: a VFS function structure, four orange rows
  pics/vptab/.style={code={
    \foreach \k in {0,...,3} \draw[white, fill=dirc] (0,-0.17*\k) rectangle (0.65,-0.17*\k-0.17);
    \foreach \k in {0,...,3} \coordinate (-r\k) at (0.65,-0.17*\k-0.085);
    \coordinate (-nw) at (0,0); \coordinate (-sw) at (0,-0.68);
    \coordinate (-ne) at (0.65,0); \coordinate (-se) at (0.65,-0.68);
  }},
}

The previous sections were about the file organization level of \Cref{sec:alloc}. This one moves up to the logical representation level: how the kernel models a file. It faces two ways: toward the user program it speaks in file descriptors and file content; toward the disk, in metadata and file content (\Cref{fig:logiclevel}).

\begin{figure}[H]
\centering
\begin{tikzpicture}
  \node[draw=usedc!60!black, fill=usedc, font=\scriptsize\bfseries, minimum width=20mm, minimum height=7mm] (user) at (0,0) {User program};
  \node[draw=bootc, fill=kspace, font=\scriptsize\bfseries, minimum width=20mm, minimum height=7mm] (kern) at (4.6,0) {Kernel};
  \pic (dk) at (8.2,0.15) {vpdisk};
  \draw[{Stealth[length=2mm]}-{Stealth[length=2mm]}, thick] (user) -- node[above, font=\scriptsize, align=center] {File descriptor and\\ file content} (kern);
  \draw[{Stealth[length=2mm]}-{Stealth[length=2mm]}, thick] (kern) -- node[above, font=\scriptsize, align=center] {Metadata and\\ file content} (dk-west |- kern);
\end{tikzpicture}
\caption{The logical representation level: the kernel between the user program and the disk (after slide 181).}
\label{fig:logiclevel}
\end{figure}

\subsection{From File Descriptor to Inode}

\begin{lstlisting}
int fd;
fd = open("/dev/zero", O_RDONLY);   // a file descriptor, if the call succeeds
\end{lstlisting}
A user program is given a file descriptor as the abstract representation of an open file. It is so abstract that it is just a number. What does the number mean?

There is a strong relationship between a process and its open files: each process has its own set of them, managed inside its \ts{} (\Cref{sec:kernel}). The chain of structures is as follows (\Cref{fig:openfile}).
\begin{itemize}
  \item The \ts{} has a pointer \code{*files} to a \code{struct files\_struct}, which manages all the open files of the process: an array \code{fd} of open files, and \code{next\_fd}, the next descriptor to hand out. By default every process starts with three open files: descriptor 0 for \code{stdin}, 1 for \code{stdout}, and 2 for \code{stderr}.
  \item When the process calls \code{open()}, the \emph{file system module} reads the file's metadata from the disk: the directory entry (with the inode number), the block addresses, and so on.
  \item The open file is represented as a \code{struct file} and added to the process. A new file descriptor, here 3, is allocated for it: the descriptor is the index of that \code{struct file} in the \code{fd} array.
  \item The \code{struct file} refers to the \emph{directory entry} of the file, which in turn contains the name and the inode of the file.
\end{itemize}

\begin{figure}[H]
\centering
\begin{tikzpicture}
  % (a) task_struct -> files_struct -> fd array; the FS module reads the disk
  \node[font=\small, anchor=east] at (-0.45,-1.2) {(a)};
  \node[vpts, minimum width=25mm, minimum height=22mm, anchor=north west] (ts) at (0,0) {};
  \node[font=\scriptsize\bfseries, text=white] at (1.25,-0.3) {struct task\_struct};
  \node[vpw] at (0.55,-0.95) {pid}; \node[vpw] at (1.6,-0.95) {counter};
  \node[vpw] (files) at (1.8,-1.75) {*files};
  \node[vpfs, minimum width=16mm, minimum height=18.5mm, anchor=north west] (fs) at (3.2,-0.35) {};
  \node[font=\scriptsize\bfseries, text=white, align=center] at (4.0,-0.75) {struct\\[-2pt] files\_struct};
  \node[vpw] at (4.0,-1.3) {next\_fd};
  \node[vpw] (fdarr) at (4.0,-1.75) {fd};
  \draw[arr] (files.east) -- (fs.west |- files);
  \foreach \k in {0,...,3} \node[vpslot] (s\k) at (5.3+0.6*\k,-1.75) {\k};
  \draw[arr] (fdarr.east) -- (s0.west);
  \foreach \k [evaluate=\k as \j using int(\k+1)] in {0,1,2} \draw[-{Stealth[length=1.3mm]}, thick] (s\k.east) -- (s\j.west);
  \node[vpmod] (mod) at (7.1,-0.6) {File system\\ module};
  \draw[arr] (mod.south -| s3) -- (s3.north);
  \node[font=\scriptsize, anchor=south west] (c1) at (2.7,0.05) {Manages all the opened files};
  \draw[vpdot] (c1.south west) ++(0.2,0) -- (files.north east);
  \node[font=\scriptsize] (c2) at (6.45,-2.55) {A list of opened files};
  \draw[vpdot] (c2.north) ++(-0.4,0) -- (s1.south);
  \begin{scope}[on background layer]
    \node[vpkern, fit={(ts)(fs)(mod)(c1)(c2)}] {};
  \end{scope}
  \pic (dk) at (8.75,-1.3) {vpdisk};
  \draw[arr] (mod.east) to[out=0, in=90] (dk-top);
  % (b) struct file -> directory entry -> the file on disk
  \begin{scope}[yshift=-3.6cm]
    \node[font=\small, anchor=east] at (-0.45,-1.2) {(b)};
    \node[vpts, fill=dirc, minimum width=26mm, minimum height=22mm, anchor=north west] (sf) at (0,0) {};
    \node[font=\small\bfseries, text=white] at (1.3,-0.35) {struct file};
    \node[vpw] (de) at (1.3,-1.1) {directory entry};
    \node[font=\scriptsize\bfseries, text=white] at (1.3,-1.7) {\dots\dots};
    \node[draw, fill=freec, rounded corners=5pt, minimum width=14mm, minimum height=16mm, anchor=north west] (dent) at (3.6,-0.5) {};
    \node[draw, font=\scriptsize, inner sep=2pt, minimum height=4mm] at (4.3,-0.82) {name};
    \node[draw, font=\scriptsize, inner sep=2pt, minimum height=4mm] (ino) at (4.3,-1.38) {inode};
    \node[font=\scriptsize\bfseries] at (4.3,-1.8) {\dots};
    \draw[arr] (de.east) -- (dent.west |- de);
    \coordinate (kr) at (6.0,0);
    \coordinate (kb) at (0,-2.55);
    \begin{scope}[on background layer]
      \node[vpkern, fit={(sf)(dent)(kr)(kb)}] {};
    \end{scope}
    \pic (doc) at (7.5,-0.45) {vpdoc};
    \pic (dk2) at (7.6,-2.0) {vpdisk};
    \draw[vpdot] (dent.east |- ino) to[out=0, in=180] (doc-west);
    \draw[vpdot] (doc-south) -- (dk2-top);
  \end{scope}
\end{tikzpicture}
\caption{An open file in the Linux kernel. (a) The \ts{} points to the \code{files\_struct}, whose \code{fd} array lists the open files; on \code{open()}, the file system module reads the metadata from the disk and the new file gets descriptor 3. (b) The \code{struct file} refers to the directory entry, which holds the name and the inode of the file (after slides 183--187).}
\label{fig:openfile}
\end{figure}

\begin{keypoint}[Scope of a File Descriptor]
  The scope of a file descriptor is restricted to its process. Different processes may use different descriptors to refer to the same file, and the same descriptor number in two processes may refer to different files (\Cref{fig:fdscope}).
\end{keypoint}

\begin{figure}[H]
\centering
\begin{tikzpicture}
  \foreach \r/\dx/\dy in {1/0/0, 2/0.45/-1.5} {
    \begin{scope}[xshift=\dx cm, yshift=\dy cm]
      \node[vpts, minimum width=13mm, minimum height=13mm] (t\r) at (0.8,-0.85) {};
      \node[vpfs, minimum width=9.5mm, minimum height=8.5mm] (f\r) at (2.35,-1.05) {};
      \foreach \k in {0,1,2} \node[vpsq] (q\r\k) at (3.45+0.55*\k,-1.05) {};
      \node[font=\scriptsize\bfseries] at (5.05,-1.05) {\dots};
      \draw[arr] (t\r.east |- f\r) -- (f\r.west);
      \draw[arr] (f\r.center) -- (q\r0.west);
      \draw[-{Stealth[length=1.3mm]}, thick] (q\r0) -- (q\r1);
      \draw[-{Stealth[length=1.3mm]}, thick] (q\r1) -- (q\r2);
    \end{scope}
  }
  \coordinate (kr) at (6.0,0); \coordinate (kb) at (0,-3.35);
  \begin{scope}[on background layer]
    \node[vpkern, fit={(t1)(t2)(f2)(kr)(kb)}] {};
  \end{scope}
  \pic (doc) at (7.4,-0.85) {vpdoc};
  \pic (dk) at (7.5,-2.45) {vpdisk};
  \draw[vpdot] (doc-south) -- (dk-top);
  \draw[arr] (q11.center) to[out=-90, in=195] (doc-west);
  \draw[arr] (q22.center) to[out=90, in=225] ($(doc-west)+(0.15,-0.5)$);
\end{tikzpicture}
\caption{Two processes refer to the same file through different file descriptors (after slide 188).}
\label{fig:fdscope}
\end{figure}

\subsection{The Virtual File System}

What is the file system module, then? It is the core of Linux's file system support: it implements how Linux deals with the disk, and it implements the \emph{virtual file system}.

\begin{definition}[Virtual File System]
  In the \emph{virtual file system} (VFS), each open file is associated with a set of file system functions: \code{open}, \code{close}, \code{read}, \code{write}, \code{lseek}, \dots{} They are just function pointers, the \emph{VFS functions}, and they are invoked by the FS-related system calls: \code{sys\_read()} calls the file's \code{read} pointer, and so on. The pointers refer to \emph{FS-specific} functions, and their values are set when the file is opened, according to where the file is: if it is on an ext4 partition, they point to the ext4 functions (\Cref{fig:vfsptr}).
\end{definition}

\begin{figure}[H]
\centering
\begin{tikzpicture}
  % the chain shared by both panels
  \newcommand{\vpchain}{%
    \node[vpts, minimum width=14mm, minimum height=14mm] (t) at (0.8,-1.2) {};
    \node[vpfs, minimum width=10mm, minimum height=9mm] (f) at (2.5,-1.45) {};
    \foreach \k in {0,1,2} \node[vpsq, minimum size=3.8mm] (q\k) at (3.7+0.65*\k,-1.45) {};
    \draw[arr] (t.east |- f) -- (f.west);
    \draw[arr] (f.center) -- (q0.west);
    \draw[-{Stealth[length=1.3mm]}, thick] (q0) -- (q1);
    \draw[-{Stealth[length=1.3mm]}, thick] (q1) -- (q2);
    \pic (ta) at (4.15,-0.25) {vptab};
    \pic (tb) at (4.65,-1.77) {vptab};
    \draw[-{Stealth[length=1.6mm]}, thick] (q0.center) -- (ta-sw);
    \draw[-{Stealth[length=1.6mm]}, thick] (q1.center) -- (tb-nw);}
  % (a) the VFS function structures
  \node[font=\small, anchor=east] at (-0.45,-1.2) {(a)};
  \vpchain
  \node[font=\scriptsize, align=center] (vl) at (7.0,-1.2) {VFS function\\ structure};
  \draw[vpdot] (vl.north west) -- (ta-r1);
  \draw[vpdot] (vl.south west) -- (tb-ne);
  \coordinate (kr) at (8.3,0); \coordinate (kb) at (0,-2.55);
  \begin{scope}[on background layer]
    \node[vpkern, fit={(t)(kr)(kb)}] {};
  \end{scope}
  % (b) the pointers refer to FS-specific functions, chosen by the file's location
  \begin{scope}[yshift=-3.7cm]
    \node[font=\small, anchor=east] at (-0.45,-1.2) {(b)};
    \vpchain
    \node[draw, cylinder, shape border rotate=90, aspect=0.3, fill=metac, minimum width=12mm, minimum height=9mm, font=\scriptsize] (cext) at (2.5,0.0) {ext4};
    \node[draw, cylinder, shape border rotate=90, aspect=0.3, fill=freec, minimum width=11mm, minimum height=7.5mm, font=\scriptsize] (cfat) at (3.35,-2.4) {FAT};
    \draw[vpdot] (q0.north) to[out=90, in=-30] (cext.south east);
    \draw[vpdot] (q1.south west) -- (cfat.north east);
    \foreach \nm/\x/\y/\c/\w in {ext4/6.6/-0.5/metac/7.5, NTFS/7.75/-0.55/vpyel/7.5, FAT/7.0/-2.05/freec/10.5} {
      \node[draw, rounded corners=4pt, fill=\c, minimum width=\w mm, minimum height=9mm] (b\nm) at (\x,\y) {};
      \node[draw, font=\tiny\bfseries, inner sep=1.5pt] at (\x,\y+0.17) {\nm};
      \node[font=\scriptsize\bfseries] at (\x,\y-0.22) {$\cdot\cdot\cdot$};
    }
    \foreach \k in {0,...,3} {
      \draw[vpdot, thin] (ta-r\k) .. controls +(0.6,0.3) and +(-0.6,0) .. ($(bext4.north west)+(0,-0.15-0.17*\k)$);
      \draw[vpdot, thin] (tb-r\k) .. controls +(0.5,0) and +(-0.5,-0.1) .. ($(bFAT.north west)+(0,-0.2-0.15*\k)$);
    }
    \coordinate (kr) at (8.3,0); \coordinate (kb) at (0,-2.6);
    \begin{scope}[on background layer]
      \node[vpkern, fit={(t)(cext)(cfat)(kr)(kb)}] {};
    \end{scope}
  \end{scope}
\end{tikzpicture}
\caption{VFS function structures. (a) Entries of the \code{fd} array point to VFS function structures. (b) The VFS function pointers refer to FS-specific functions; which ones is determined by the location of the opened file, an ext4 or a FAT partition (after slides 190--193).}
\label{fig:vfsptr}
\end{figure}

\begin{figure}[H]
\centering
\begin{tikzpicture}[hdr/.style={draw, font=\scriptsize\bfseries, inner sep=2pt},
                    fsb/.style={draw, rounded corners=10pt, minimum width=29mm, minimum height=26mm},
                    fat/.style={single arrow, draw, minimum width=8mm, single arrow head extend=2mm, inner sep=1.5mm}]
  \newcommand{\vpfat}[3]{%
    \path let \p1=($(#2)-(#1)$), \n1={atan2(\y1,\x1)}, \n2={veclen(\x1,\y1)} in
      node[fat, fill=#3, minimum height=\n2, rotate=\n1, anchor=tail] at (#1) {};}
  \node[draw, fill=black!15, minimum width=48mm, minimum height=17mm] (sys) at (0,0) {};
  \node[hdr, fill=black!15] at (0,0.52) {System calls};
  \node[font=\scriptsize\ttfamily, align=center] at (0,-0.22) {sys\_open(...); sys\_close(...);\\ sys\_read(...); sys\_write(...);\\ sys\_lseek(...);};
  \node[draw, fill=dirc, minimum width=48mm, minimum height=17mm] (vfs) at (6.6,0) {};
  \node[hdr, fill=dirc] at (6.6,0.52) {VFS functions};
  \node[font=\scriptsize\ttfamily, align=center] at (6.6,-0.22) {open(...); close(...);\\ read(...); write(...);\\ lseek(...);};
  \node[single arrow, fill=freec, minimum height=15mm, minimum width=9mm, single arrow head extend=1.5mm, inner sep=1pt, font=\scriptsize\bfseries] at (3.3,0) {Invoke};
  \foreach \nm/\pre/\x/\c in {ext4/ext4/0.1/metac, FAT/fat/3.5/freec, NTFS/ntfs/6.9/vpyel} {
    \node[fsb, fill=\c] (b\pre) at (\x,-3.7) {};
    \node[hdr] at (\x,-2.75) {\nm{} FS functions};
    \node[font=\scriptsize\ttfamily, align=left] at (\x,-3.9) {\pre\_open(...)\\ \pre\_close(...)\\ \pre\_read(...)\\ \pre\_write(...)\\ \pre\_lseek(...)};
  }
  \vpfat{4.4,-0.95}{1.85,-2.4}{metac}
  \vpfat{5.6,-1.0}{4.45,-2.4}{freec}
  \vpfat{6.9,-0.95}{6.9,-2.35}{vpyel}
  \node[font=\scriptsize, align=left, anchor=east] at (2.7,-1.45) {Invoking VFS functions means\\ invoking the ext4 functions.};
  \node[font=\scriptsize, anchor=west] at (7.35,-1.55) {Pointer assignments};
\end{tikzpicture}
\caption{The VFS: system calls invoke the VFS function pointers of the open file, which were assigned the functions of the file system the file lives on. For a file on ext4, invoking the VFS functions means invoking the ext4 functions (after slides 191 and 194).}
\label{fig:vfs}
\end{figure}

This is the function-pointer technique of Lab~1 (\code{nyuc}) applied at kernel scale: a table of pointers is an interface, and each file system is an implementation of it.

\emph{Pros.} The VFS lets the OS support multiple file systems at the same time, although the \code{open()} for ext4 is a lot different from the \code{open()} for FAT32. And it is easy to support a new file system in the future: all one has to do is implement the FS functions specified by the VFS structure. This is why Linux supports so many file systems.

\emph{Cons.} The VFS cannot make use of special functions of a specific file system. Suppose some FS offers an operation that combines \code{open()} and \code{read()} for small, read-only files. The VFS structure has no slot for it, so no program can reach it through the system call interface.
